% !TeX program = xelatex
% UTF-8; compile twice with XeLaTeX. No external images or data files needed.
\documentclass[a4paper,landscape,10pt]{article}
\usepackage[margin=14mm,headheight=14pt,headsep=5mm,footskip=7mm]{geometry}
\usepackage{fontspec}
\setmainfont{DejaVu Sans}
\newfontfamily\ipafont{DejaVu Sans}
\usepackage[table]{xcolor}
\usepackage{array,longtable,booktabs,fancyhdr}
\definecolor{ink}{HTML}{18384D}
\definecolor{band}{HTML}{E7EFF4}
\pagestyle{fancy}\fancyhf{}
\fancyhead[L]{\small\color{ink}Parlez Vous QSO?! · Latviski / Latvian}
\fancyhead[R]{\small Englische Phrase · Lettische Übersetzung · IPA}
\fancyfoot[L]{\scriptsize Vorlage: Latvian v 1.0 · Credits: Agris -- YL2VW}
\fancyfoot[R]{\small\thepage}
\setlength{\parindent}{0pt}
\setlength{\tabcolsep}{3mm}
\renewcommand{\arraystretch}{1.25}
\newcommand{\ipa}[1]{{\ipafont[#1]}}
\begin{document}
{\LARGE\bfseries\color{ink}Parlez Vous QSO?!}\hfill{\large Latviski / Latvian}

\vspace{2mm}
{\small QSO-Sprachführer mit lettischer IPA-Lautschrift. Die Lautschrift bezieht sich auf die lettische Spalte.}

\vspace{2mm}
{\footnotesize \textbf{Leseschlüssel:} \ipa{ˈ} Hauptbetonung, \ipa{ˌ} Nebenbetonung, \ipa{ː} langer Vokal;
\ipa{c} = ķ, \ipa{ɟ} = ģ, \ipa{ɲ} = ņ, \ipa{ʎ} = ļ, \ipa{ts} = c, \ipa{tʃ} = č, \ipa{ʃ} = š.
Grundlage: Standardlettisch, ohne Silbentöne. Schrägstriche trennen Alternativen; \ipa{…} steht für Platzhalter.}

\vspace{1mm}
\fontsize{9}{11.8}\selectfont
\begin{longtable}{@{}>{\raggedright\arraybackslash}p{13mm}>{\raggedright\arraybackslash}p{73mm}>{\raggedright\arraybackslash}p{75mm}>{\raggedright\arraybackslash}p{90mm}@{}}
\toprule
\textbf{Ref.} & \textbf{English phrase} & \textbf{Latviski} & \textbf{Lautschrift (IPA)}\\
\midrule
\endfirsthead
\toprule
\textbf{Ref.} & \textbf{English phrase} & \textbf{Latviski} & \textbf{Lautschrift (IPA)}\\
\midrule
\endhead
\midrule
\multicolumn{4}{r}{\scriptsize Fortsetzung auf der nächsten Seite}\\
\endfoot
\bottomrule
\endlastfoot
\rowcolor{band} 1.0 & \textbf{Opening Greetings} & \multicolumn{2}{l}{\textbf{Sarunas uzsākšana}} \\*
1.1 & I am very pleased to meet you. & Man prieks jūs satikt. & \ipa{mɑn prieks juːs ˈsɑtikt} \\
1.2 & Thank you for your call & Paldies par izsaukumu. & \ipa{ˈpɑldies pɑr ˈizsɑukumu} \\
1.3 & Good morning & Labrīt. & \ipa{ˈlɑbriːt} \\
1.31 & Good day & Labdien. & \ipa{ˈlɑbdien} \\
1.32 & Good afternoon & Labdien. & \ipa{ˈlɑbdien} \\
1.33 & Good evening & Labvakar. & \ipa{ˈlɑbvɑkɑr} \\
1.4 & I know only a little Latvian but I enjoy QSOs in Latvian & Es māku tikai nedaudz latviski, bet man patīk veidot QSO latviski. & \ipa{æs ˈmɑːku ˈtikai ˈnedɑudz ˈlɑtviski bæt mɑn ˈpɑtiːk ˈveiduɔt kuː es ɔː ˈlɑtviski} \\
\rowcolor{band} 2.0 & \textbf{Operator} & \multicolumn{2}{l}{\textbf{Operators}} \\*
2.1 & My name is Peter. & Mans vārds ir Pēteris. & \ipa{mɑns vɑːrts ir ˈpeːteris} \\
2.2 & I have been licensed since 1960. & Esmu licencēts kopš 1960 gada. & \ipa{ˈæsmu ˈlitsentseːts kuɔpʃ ˈtuːkstuɔʃ ˈdeviɲi ˈsimti ˈseʒdesmitɑː ˈɡɑdɑ} \\
2.3 & I have been licensed for 43 years. & Esmu licencēts 43 gadus. & \ipa{ˈæsmu ˈlitsentseːts ˈtʃetrdesmit triːs ˈɡɑdus} \\
2.4 & I live in Sittingbourne. & Es dzīvoju Sittingbourne. & \ipa{æs ˈdziːvuɔju ˈsɪtɪŋbɔːn} \\
2.5 & Sittingbourne is a village. & Sittingbourne ir ciemats. & \ipa{ˈsɪtɪŋbɔːn ir ˈtsiemɑts} \\
2.51 & Sittingbourne is a town. & Sittingbourne ir mazpilsēta. & \ipa{ˈsɪtɪŋbɔːn ir ˈmɑspilseːtɑ} \\
2.52 & Sittingbourne is a city. & Sittingbourne ir lielpilsēta. & \ipa{ˈsɪtɪŋbɔːn ir ˈlielpilseːtɑ} \\
2.6 & Sittingbourne is between London and Dover. & Sittingbourne ir starp Londonu un Doveru. & \ipa{ˈsɪtɪŋbɔːn ir stɑrp ˈlɔndɔnu un ˈdəʊvəru} \\
2.7 & I live in the South of England. & Es dzīvoju Anglijas dienvidos. & \ipa{æs ˈdziːvuɔju ˈɑŋɡlijɑs ˈdienviduɔs} \\
2.71 & I live in the North of England. & Es dzīvoju Anglijas ziemeļos. & \ipa{æs ˈdziːvuɔju ˈɑŋɡlijɑs ˈziemeʎuɔs} \\
2.72 & I live in the East of England. & Es dzīvoju Anglijas austrumos. & \ipa{æs ˈdziːvuɔju ˈɑŋɡlijɑs ˈɑustrumuɔs} \\
2.73 & I live in West of England. & Es dzīvoju Anglijas rietumos. & \ipa{æs ˈdziːvuɔju ˈɑŋɡlijɑs ˈrietumuɔs} \\
2.74 & I live in South East England. & Es dzīvoju Anglijas dienvidaustrumos. & \ipa{æs ˈdziːvuɔju ˈɑŋɡlijɑs ˈdienvidˌɑustrumuɔs} \\
2.75 & I live in South West England. & Es dzīvoju Anglijas dienvidrietumos. & \ipa{æs ˈdziːvuɔju ˈɑŋɡlijɑs ˈdienvidˌrietumuɔs} \\
2.76 & I live in North East England. & Es dzīvoju Anglijas ziemeļaustrumos. & \ipa{æs ˈdziːvuɔju ˈɑŋɡlijɑs ˈziemeʎˌɑustrumuɔs} \\
2.77 & I live in North West England. & Es dzīvoju Anglijas ziemeļrietumos. & \ipa{æs ˈdziːvuɔju ˈɑŋɡlijɑs ˈziemeʎˌrietumuɔs} \\
2.78 & I live in central England. & Es dzīvoju centrālajā Anglijā. & \ipa{æs ˈdziːvuɔju ˈtsentrɑːlɑjɑː ˈɑŋɡlijɑː} \\
2.8 & I live in the country. & Es dzīvoju laukos. & \ipa{æs ˈdziːvuɔju ˈlɑukuɔs} \\
2.81 & I live by the sea. & Es dzīvoju pie jūras. & \ipa{æs ˈdziːvuɔju pie ˈjuːrɑs} \\
2.82 & I live near mountains. & Es dzīvoju blakus kalniem. & \ipa{æs ˈdziːvuɔju ˈblɑkus ˈkɑlniem} \\
2.9 & Sittingbourne is 50Km from London. & Sittingbourne ir 50 km no Londonas. & \ipa{ˈsɪtɪŋbɔːn ir ˈpiedzdesmit ˈkilɔmetru nuɔ ˈlɔndɔnɑs} \\
2.11 & I am retired. & Es esmu pensijā. & \ipa{æs ˈæsmu ˈpensijɑː} \\
2.12 & My home call sign is M0ECS. & Mans mājas izsaukuma signāls ir M0ECS. & \ipa{mɑns ˈmɑːjɑs ˈizsɑukumɑ ˈsiɡnɑːls ir em ˈnulːe eː tseː es} \\
\rowcolor{band} 3.0 & \textbf{Station} & \multicolumn{2}{l}{\textbf{Stacija}} \\*
3.1 & My radio is an Icom 706. & Mans radio ir Icom 706. & \ipa{mɑns ˈrɑdiɔ ir ˈaɪkɒm ˈseptiɲi ˈsimti ˈseʃi} \\
3.2 & My station is all homebrew (homemade). & Mans radio ir paštaisīts. & \ipa{mɑns ˈrɑdiɔ ir ˈpɑʃtɑisiːts} \\
3.3 & I am running 5 Watts. & Es strādāju ar 5 W jaudu. & \ipa{æs ˈstrɑːdɑːju ɑr ˈpietsu ˈvɑtu ˈjɑudu} \\
3.4 & My antenna is a long wire. & Mana antena ir "long wire". & \ipa{ˈmɑnɑ ˈɑntenɑ ir ˈlɒŋ ˈwaɪə} \\
3.41 & My antenna is a dipole. & Mana antena ir dipols. & \ipa{ˈmɑnɑ ˈɑntenɑ ir ˈdipɔls} \\
3.42 & My antenna is a G5RV. & Mana antena ir G5RV. & \ipa{ˈmɑnɑ ˈɑntenɑ ir ɡeː pietsi er veː} \\
3.43 & My antenna is a doublet. & Mana antena ir doublet. & \ipa{ˈmɑnɑ ˈɑntenɑ ir ˈdʌblɪt} \\
3.44 & My antenna is a delta loop. & Mana antena ir trijstūris. & \ipa{ˈmɑnɑ ˈɑntenɑ ir ˈtrijstuːris} \\
3.45 & My antenna is a vertical. & Mana ir vertikāla antena. & \ipa{ˈmɑnɑ ir ˈvertikɑːlɑ ˈɑntenɑ} \\
3.46 & My antenna is a log periodic. & Man ir logoperiodiska antena. & \ipa{mɑn ir ˈlɔɡɔperiɔdiskɑ ˈɑntenɑ} \\
3.47 & My antenna is a Yagi beam. & Man ir Yagi antena. & \ipa{mɑn ir ˈjɑɡi ˈɑntenɑ} \\
3.48 & My antenna is a 5 element beam. & Man ir 5 elementu virziendarbības antena. & \ipa{mɑn ir ˈpietsu ˈelementu ˈvirzienˌdɑrbiːbɑs ˈɑntenɑ} \\
\rowcolor{band} 4.0 & \textbf{Conditions} & \multicolumn{2}{l}{\textbf{Apstākļi}} \\*
4.1 & Thank you for my 599 report. & Paldies par 599. & \ipa{ˈpɑldies pɑr pietsi ˈdeviɲi ˈdeviɲi} \\
4.2 & The band conditions are good. & Izplatīšanās diapazonā ir laba. & \ipa{ˈizplɑtiːʃɑnɑːs ˈdiɑpɑzɔnɑː ir ˈlɑbɑ} \\
4.3 & The band conditions are bad. & Izplatīšanās diapazonā ir slikta. & \ipa{ˈizplɑtiːʃɑnɑːs ˈdiɑpɑzɔnɑː ir ˈsliktɑ} \\
4.4 & There is strong interference here. & Ir spēcīgi traucējumi. & \ipa{ir ˈspeːtsiːɡi ˈtrɑutseːjumi} \\
4.5 & Can you change to the 80m band? & Vai jūs varat pāriet uz 80 m diapazonu? & \ipa{vai juːs ˈvɑrɑt ˈpɑːriet uz ˈɑstuɔɲdesmit ˈmetru ˈdiɑpɑzɔnu} \\
4.6 & Please change frequency to 3.558 MHz. & Lūdzu nomainiet frekvenci uz 3,558 MHz. & \ipa{ˈluːdzu ˈnuɔmainiet ˈfrekventsi uz triːs ˈkɔmɑts pietsi pietsi ˈɑstuɔɲi ˈmeɡɑhertsi} \\
\rowcolor{band} 5.0 & \textbf{General Queries and Comments} & \multicolumn{2}{l}{\textbf{Vispārīgi jautājumi un komentāri}} \\*
5.1 & What does <\#\#\#\#\#> mean please? & Ko nozīmē <\#\#\#\#\#>? & \ipa{kuɔ ˈnuɔziːmeː …} \\
5.2 & Please wait a moment, I need to check my Latvian dictionary. & Lūdzu uzgaidiet, es paskatīšos latviešu valodas vārdnīcā. & \ipa{ˈluːdzu ˈuzɡaidiet æs ˈpɑskɑtiːʃuɔs ˈlɑtvieʃu ˈvɑluɔdɑs ˈvɑːrdniːtsɑː} \\
5.3 & How do you say, <\#\#\#\#\#> in Latvian? & Kā latviski pateikt <\#\#\#\#\#>? & \ipa{kɑː ˈlɑtviski ˈpɑteikt …} \\
5.5 & My favourite modes of operation are CW / RTTY / PSK / MT63 / Olivia / Hellschreiber / Digital Modes / and SSB & Mani iecienītākie darba veidi ir CW / RTTY / PSK / MT63 / Olivia / Hellschreiber / ciparu veidi / un SSB. & \ipa{ˈmɑni ˈietsieniːtɑːkie ˈdɑrbɑ ˈveidi ir tseː ˈdubult veː / er teː teː ˈipsilɔn / peː es kɑː / em teː ˈseʒdesmit triːs / oˈliviɑ / ˈhɛlʃraɪbɐ / ˈtsipɑru ˈveidi / un es es beː} \\
5.6 & My best DX is Australia. & Mans labākais DX ir Austrālija. & \ipa{mɑns ˈlɑbɑːkais deː iks ir ˈɑustrɑːlijɑ} \\
5.7 & I enjoy working contests. & Man patīk piedalīties sacensībās. & \ipa{mɑn ˈpɑtiːk ˈpiedɑliːties ˈsɑtsensiːbɑːs} \\
5.8 & I like working QRP. & Man patīk strādāt QRP. & \ipa{mɑn ˈpɑtiːk ˈstrɑːdɑːt kuː er peː} \\
\rowcolor{band} 6.0 & \textbf{The Weather} & \multicolumn{2}{l}{\textbf{Laika apstākļi}} \\*
6.1 & It is sunny. & Ir saulains. & \ipa{ir ˈsɑulɑins} \\
6.2 & It is windy. & Ir vējains. & \ipa{ir ˈveːjɑins} \\
6.3 & It is raining. & Līst lietus. & \ipa{liːst ˈlietus} \\
6.4 & It is snowing. & Snieg sniegs. & \ipa{sniek snieks} \\
6.5 & It is very hot. & Ir ļoti karsts. & \ipa{ir ˈʎuɔti kɑrsts} \\
6.6 & It is very cold. & Ir ļoti auksts. & \ipa{ir ˈʎuɔti ɑuksts} \\
6.7 & The temperature is 8C. & Temperatūra ir 8 °C. & \ipa{ˈtemperɑtuːrɑ ir ˈɑstuɔɲi ˈɡrɑːdi peːts ˈtselsijɑ} \\
6.8 & It is overcast. & Ir apmācies. & \ipa{ir ˈɑpmɑːtsies} \\
6.9 & It is foggy. & Ir migla. & \ipa{ir ˈmiɡlɑ} \\
6.11 & It has been raining all day. & Visu dienu līst. & \ipa{ˈvisu ˈdienu liːst} \\
6.12 & There is 20cm of snow here. & Te ir 20 cm sniega. & \ipa{te ir ˈdivdesmit ˈtsentimetru ˈsnieɡɑ} \\
\rowcolor{band} 8.0 & \textbf{QSL} & \multicolumn{2}{l}{\textbf{QSL}} \\*
8.1 & My QSL card is sure via the bureau & QSL kartiņu noteikti nosūtīšu caur biroju. & \ipa{kuː es el ˈkɑrtiɲu ˈnuɔteikti ˈnuɔsuːtiːʃu tsɑur ˈbirɔju} \\
8.2 & QSL via eQSL & QSL caur eQSL. & \ipa{kuː es el tsɑur iː kuː es el} \\
8.3 & No eQSL please & Lūdzu nesūtiet eQSL. & \ipa{ˈluːdzu ˈnesuːtiet iː kuː es el} \\
8.4 & I will send you my card direct. & Es nosūtīšu jums savu kartiņu pa tiešo. & \ipa{æs ˈnuɔsuːtiːʃu jums ˈsɑvu ˈkɑrtiɲu pɑ ˈtieʃuɔ} \\
8.5 & What is your address? & Kāda ir jūsu adrese? & \ipa{ˈkɑːdɑ ir ˈjuːsu ˈɑdrese} \\
8.6 & Please send your QSL card direct. & Lūdzu sūtiet jūsu QSL pa tiešo. & \ipa{ˈluːdzu ˈsuːtiet ˈjuːsu kuː es el pɑ ˈtieʃuɔ} \\
8.7 & My address is... & Mana adrese ir… & \ipa{ˈmɑnɑ ˈɑdrese ir …} \\
8.8 & My address is shown on QSL.com / EHam.net & Mana adrese ir atrodama QSL.com / EHam.net. & \ipa{ˈmɑnɑ ˈɑdrese ir ˈɑtruɔdɑmɑ kuː es el ˈpuŋkts kɔm / iː hæm ˈpuŋkts net} \\
8.9 & Sorry I do not send QSL cards. & Piedodiet, bet es nesūtu QSL. & \ipa{ˈpieduɔdiet bæt æs ˈnesuːtu kuː es el} \\
8.11 & Please send QSL card via G1JQH. & Lūdzu sūtiet QSL kartiņu caur G1JQH. & \ipa{ˈluːdzu ˈsuːtiet kuː es el ˈkɑrtiɲu tsɑur ɡeː viens jeː kuː hɑː} \\
8.12 & My WACRAL / GQRP / DXCC / FISTS number is 1234. & Mans WACRAL / GQRP / DXCC / FISTS numurs ir 1234. & \ipa{mɑns ˈwakral / ɡeː kuː er peː / deː iks tseː tseː / fɪsts ˈnumurs ir ˈtuːkstuɔʃ ˈdivi ˈsimti ˈtriːsdesmit ˈtʃetri} \\
\rowcolor{band} 9.0 & \textbf{Final Greetings} & \multicolumn{2}{l}{\textbf{Atvadīšanās}} \\*
9.1 & Thank you. & Paldies. & \ipa{ˈpɑldies} \\
9.2 & Thank you for a very enjoyable contact. & Paldies par ļoti patīkamu sakaru. & \ipa{ˈpɑldies pɑr ˈʎuɔti ˈpɑtiːkɑmu ˈsɑkɑru} \\
9.3 & I enjoyed our chat. & Bija patīkami papļāpāt. & \ipa{ˈbijɑ ˈpɑtiːkɑmi ˈpɑpʎɑːpɑːt} \\
9.4 & I enjoyed our QSO in Latvian & Man prieks par šo QSO latviešu valodā. & \ipa{mɑn prieks pɑr ʃuɔ kuː es ɔː ˈlɑtvieʃu ˈvɑluɔdɑː} \\
9.5 & It has been great fun. & Bija forši. & \ipa{ˈbijɑ ˈfɔrʃi} \\
9.6 & If you would like a regular contact with me, you can contact me via email. & Ja jūs vēlaties sazināties ar mani, lūdzu izmantojiet e-pastu. & \ipa{jɑ juːs ˈveːlɑties ˈsɑzinɑːties ɑr ˈmɑni ˈluːdzu ˈizmɑntuɔjiet ˈeːˌpɑstu} \\
9.7 & My email address is... & Mana e-pasta adrese ir… & \ipa{ˈmɑnɑ ˈeːˌpɑstɑ ˈɑdrese ir …} \\
9.8 & My email address is shown on QSL.com / EHam.net & Mana e-pasta adrese ir atrodama QSL.com / EHam.net. & \ipa{ˈmɑnɑ ˈeːˌpɑstɑ ˈɑdrese ir ˈɑtruɔdɑmɑ kuː es el ˈpuŋkts kɔm / iː hæm ˈpuŋkts net} \\
9.9 & I look forward to meeting you again. & Būšu priecīgs jūs atkal satikt. & \ipa{ˈbuːʃu ˈprietsiːɡs juːs ˈɑtkɑl ˈsɑtikt} \\
9.11 & My best wishes to you and your family. & Vislabākie novēlējumi jums un jūsu ģimenei. & \ipa{ˈvislɑbɑːkie ˈnuɔveːleːjumi jums un ˈjuːsu ˈɟimenei} \\
9.12 & Best wishes to your wife. & Labākie novēlējumi jūsu sievai. & \ipa{ˈlɑbɑːkie ˈnuɔveːleːjumi ˈjuːsu ˈsievɑi} \\
9.121 & Best wishes to your husband. & Labākie novēlējumi jūsu vīram. & \ipa{ˈlɑbɑːkie ˈnuɔveːleːjumi ˈjuːsu ˈviːrɑm} \\
9.13 & Good luck. & Laimīgi. & \ipa{ˈlaimiːɡi} \\
9.14 & Good health. & Labu veselību. & \ipa{ˈlɑbu ˈveseliːbu} \\
9.15 & Goodbye & Uzredzēšanos. & \ipa{ˈuzredzeːʃɑnuɔs} \\
9.16 & I wish you an enjoyable afternoon. & Lai jums jauka pēcpusdiena. & \ipa{lai jums ˈjɑukɑ ˈpeːtsˌpusdienɑ} \\
9.161 & I wish you an enjoyable evening. & Lai jums jauks vakars. & \ipa{lai jums jɑuks ˈvɑkɑrs} \\
9.162 & I wish you an enjoyable weekend. & Lai jums jauka nedēļas nogale. & \ipa{lai jums ˈjɑukɑ ˈnedeːʎɑs ˈnuɔɡɑle} \\
9.17 & Enjoy your holiday. & Lai jums jaukas brīvdienas. & \ipa{lai jums ˈjɑukɑs ˈbriːvˌdienɑs} \\
9.18 & Happy New Year. & Laimīgu jauno gadu. & \ipa{ˈlaimiːɡu ˈjɑunuɔ ˈɡɑdu} \\
\rowcolor{band} 10.0 & \textbf{Country} & \multicolumn{2}{l}{\textbf{Valsts}} \\*
10.1 & Albania & Albānija & \ipa{ˈɑlbɑːnijɑ} \\
10.2 & Austria & Austrija & \ipa{ˈɑustrijɑ} \\
10.3 & Belgium & Beļģija & \ipa{ˈbeʎɟijɑ} \\
10.4 & Belorussia & Baltkrievija & \ipa{ˈbɑltˌkrievijɑ} \\
10.5 & Bosnia \& Herzegovina & Bosnija un Hercegovina & \ipa{ˈbɔsnijɑ un ˈhertseɡɔvinɑ} \\
10.6 & Bulgaria & Bulgārija & \ipa{ˈbulɡɑːrijɑ} \\
10.7 & Croatia & Horvātija & \ipa{ˈhɔrvɑːtijɑ} \\
10.8 & Czech Republic & Čehija & \ipa{ˈtʃehijɑ} \\
10.9 & Denmark & Dānija & \ipa{ˈdɑːnijɑ} \\
10.11 & Eire & Īrija & \ipa{ˈiːrijɑ} \\
10.12 & England & Anglija & \ipa{ˈɑŋɡlijɑ} \\
10.13 & Estonia & Igaunija & \ipa{ˈiɡɑunijɑ} \\
10.14 & Finland & Somija & \ipa{ˈsuɔmijɑ} \\
10.15 & France & Francija & \ipa{ˈfrɑntsijɑ} \\
10.16 & Germany & Vācija & \ipa{ˈvɑːtsijɑ} \\
10.17 & Gibraltar & Gibraltārs & \ipa{ˈɡibrɑltɑːrs} \\
10.18 & Greece & Grieķija & \ipa{ˈɡriecijɑ} \\
10.19 & Holland & Holande & \ipa{ˈhɔlɑnde} \\
10.21 & Hungary & Ungārija & \ipa{ˈuŋɡɑːrijɑ} \\
10.22 & Iceland & Islande & \ipa{ˈislɑnde} \\
10.23 & Northern Ireland & Ziemeļīrija & \ipa{ˈziemeʎˌiːrijɑ} \\
10.24 & Isle of Man & Menas sala & \ipa{ˈmenɑs ˈsɑlɑ} \\
10.25 & Italy & Itālija & \ipa{ˈitɑːlijɑ} \\
10.26 & Latvia & Latvija & \ipa{ˈlɑtvijɑ} \\
10.27 & Liechtenstein & Lihtenšteina & \ipa{ˈlihtenʃteinɑ} \\
10.28 & Lithuania & Lietuva & \ipa{ˈlietuvɑ} \\
10.29 & Luxembourg & Luksemburga & \ipa{ˈluksemburɡɑ} \\
10.31 & Macedonia & Maķedonija & \ipa{ˈmɑcedɔnijɑ} \\
10.32 & Moldavia & Moldāvija & \ipa{ˈmɔldɑːvijɑ} \\
10.33 & Netherlands & Nīderlande & \ipa{ˈniːderlɑnde} \\
10.34 & Norway & Norvēģija & \ipa{ˈnɔrveːɟijɑ} \\
10.35 & Poland & Polija & \ipa{ˈpuɔlijɑ} \\
10.36 & Portugal & Portugāle & \ipa{ˈpɔrtuɡɑːle} \\
10.37 & Romania & Rumānija & \ipa{ˈrumɑːnijɑ} \\
10.38 & Russian Federation & Krievijas Federācija & \ipa{ˈkrievijɑs ˈfederɑːtsijɑ} \\
10.39 & Scotland & Skotija & \ipa{ˈskɔtijɑ} \\
10.41 & Slovakia & Slovākija & \ipa{ˈslɔvɑːkijɑ} \\
10.42 & Slovenia & Slovēnija & \ipa{ˈslɔveːnijɑ} \\
10.43 & Spain & Spānija & \ipa{ˈspɑːnijɑ} \\
10.44 & Svalbard & Svalbarda & \ipa{ˈsvɑlbɑrdɑ} \\
10.45 & Sweden & Zviedrija & \ipa{ˈzviedrijɑ} \\
10.46 & Switzerland & Šveice & \ipa{ˈʃveitse} \\
10.47 & Turkey & Turcija & \ipa{ˈturtsijɑ} \\
10.48 & Ukraine & Ukraina & \ipa{ˈukrainɑ} \\
10.49 & Wales & Velsa & \ipa{ˈvelsɑ} \\
\end{longtable}
\normalsize
\section*{Hinweise zur Übertragung und Aussprache}
\fontsize{8.5}{10}\selectfont\raggedright
\textbf{Vorlage und Umfang.} Grundlage ist „Latvian v 1.0.doc“, Credit: Agris, YL2VW.
Alle 142 Phrasen und Ländernamen sind enthalten. Die Gliederung und ursprünglichen Nummernsprünge
bleiben erhalten, einschließlich des fehlenden Abschnitts 7. Die zwei leeren zusätzlichen Spalten der
Word-Vorlage wurden nicht übernommen; die vierte Spalte enthält die Lautschrift.

\smallskip
\textbf{Korrekturen.} „licenzēts“ wurde zu „licencēts“ (2.2/2.3), „Gada“ zu „gada“ (2.2).
Die Himmelsrichtungen dienvidaustrumos, dienvidrietumos, ziemeļaustrumos und ziemeļrietumos
stehen ohne Bindestrich (2.74--2.77). „centrāl Anglijā“ wurde zu „centrālajā Anglijā“ (2.78).
Icom706 steht als Icom 706 (3.1), 80m als 80 m (4.5), 3.558 MHz mit Dezimalkomma als 3,558 MHz (4.6).
MT35 wurde gemäß der englischen Spalte zu MT63 (5.5); „Latviešu“, „Ciparu“ und „Jauno“
stehen klein (5.2/5.5/9.18). „izmanojiet“ ist „izmantojiet“ (9.6).
„Hercogovina“ wurde zu „Hercegovina“ (10.5), „Gibraltāra“ zum Ländernamen „Gibraltārs“ (10.17).
Temperatur und Auslassungen wurden typografisch vereinheitlicht.

\smallskip
\textbf{Quellenbesonderheiten.} Sittingbourne bleibt in der lettischen Spalte in der englischen Form,
auch ohne lettische Ortsendung in 2.4. 3.2 beschreibt das selbstgebaute Radio statt die gesamte Station.
3.44 nennt ein Dreieck (trijstūris), statt ausdrücklich eine Delta-Loop-Antenne.
4.1 nennt die Ziffern 599 ohne das Wort Rapport.
9.6 beschreibt eine Kontaktaufnahme per E-Mail ohne ausdrücklich einen regelmäßigen Kontakt zu nennen.
Die IPA-Spalte folgt diesen lettischen Formulierungen. Historische Ländernamen und Adressen bleiben erhalten.

\smallskip
\textbf{Aussprachebasis.} Die manuelle IPA-Spalte ist eine breite Lesehilfe für Standardlettisch.
\ipa{ˈ} bezeichnet die Hauptbetonung, die meist auf der ersten Silbe liegt;
\ipa{ˌ} ausgewählte Nebenbetonungen in zusammengesetzten Wörtern.
\ipa{ː} kennzeichnet lange Vokale, entsprechend ā, ē, ī und ū.
Die lettischen Silbenintonationen, etwa gedehnt, fallend oder gebrochen, werden nicht eingetragen;
Satzmelodie, feine Konsonantenlängen und alle Assimilationen sind ebenfalls nicht vollständig kodiert.

\smallskip
\textbf{Vokale und Konsonanten.} Ķ und ģ stehen als palatale Verschlusslaute \ipa{c, ɟ},
ņ als \ipa{ɲ}, ļ als \ipa{ʎ}. Geschriebenes c ist \ipa{ts}, č ist \ipa{tʃ}; dz ist \ipa{dz}.
Der Buchstabe o steht je nach Wort für \ipa{uɔ} oder \ipa{ɔ}; der Diphthong tritt etwa in dzīvoju und no auf.
Ie und andere Diphthonge stehen vereinfacht als Vokalfolgen.
Enges und breites e werden hier vereinfacht mit \ipa{e} und in ausgewählten häufigen Wörtern \ipa{æ}
dargestellt, etwa es und bet. Die vollständige Verteilung der e/ē-Qualitäten ist nicht kodiert.
R steht als \ipa{r}. Einige stimmlose Varianten vor Konsonanten sind berücksichtigt,
etwa vārds mit \ipa{t} und mazpilsēta mit \ipa{s}.

\smallskip
\textbf{Namen und Fachwörter.} Pēteris ist lettisch gelesen; Sittingbourne, Icom, long wire und doublet
verwenden englische Lesarten. Doveru kombiniert eine englische Namenslesart mit der lettischen Endung.
Hellschreiber steht deutsch. Yagi und die Ländernamen stehen überwiegend mit lettischer Aussprache,
einschließlich Lihtenšteina und Svalbarda. Namensvarianten sind möglich.
Rufzeichen und Kürzel verwenden hier Q „kū“, S „es“, L „el“, G „gē“, R „er“, C „cē“, J „jē“, H „hā“;
W als „dubult vē“, Y als „ipsilon“. WACRAL und FISTS werden als Wörter gelesen.
E in eQSL/eHam ist \ipa{iː}; e-pasts verwendet das lettische \ipa{eː}.
Webpunkte werden als „punkts“ gesprochen, com/net als Wörter.

\smallskip
\textbf{Zahlen.} 1960 ist nach kopš als „tūkstoš deviņi simti sešdesmitā“ gelesen,
43 als „četrdesmit trīs“, 50 als „piecdesmit“, 706 als „septiņi simti seši“.
5 W wird im Satz mit „piecu vatu“, 80 m mit „astoņdesmit metru“,
20 cm mit „divdesmit centimetru“ aufgelöst.
599 steht ziffernweise als „pieci deviņi deviņi“, MT63 mit „sešdesmit trīs“;
3,558 MHz als „trīs komats pieci pieci astoņi megaherci“.
1234 = „tūkstoš divi simti trīsdesmit četri“.
8 °C wird als „astoņi grādi pēc Celsija“ gelesen. \ipa{…} steht für offene Angaben und Platzhalter.

\smallskip
\textbf{Fachhinweis.} A. Krišjānis Kariņš: „The Prosodic Structure of Latvian“, Dissertation,
University of Pennsylvania, 1996, besonders Kapitel 2 zum Lautsystem sowie die Kapitel zu Betonung und
Silbenintonation. Öffentliches Exemplar: languagelog.ldc.upenn.edu/myl/Karins1996.pdf.
Die Konventionen wurden für diese vereinfachte Lesehilfe angepasst.
\end{document}
